% Part: lambda-calculus
% Chapter: introduction
% Section: beta

\documentclass[../../../include/open-logic-section]{subfiles}

\begin{document}

\olfileid[hi]{lam}{int}{bet}
\olsection{$\beta$-अपचयन}

$(\lambd[m][(\lambd[y][y]) m])$ को देखकर यह अनुमान स्वाभाविक है कि उसका
$\lambd[m][m]$ से कोई संबंध है; अर्थात् दूसरा पद पहले को ``सरल करने'' का
परिणाम होना चाहिए। \emph{$\beta$-अपचयन} की धारणा इस अंतर्बोध को सटीक रूप
देती है।

\begin{defn}[$\beta$-संकुचन, $\bredone$] \ollabel{defn:betacontr}
  \emph{$\beta$-संकुचन} ($\bredone$) पदों पर वह सबसे छोटा संगत संबंध है
  जो निम्न शर्त संतुष्ट करता है:
  \[
    (\lambd[x][N])Q \bredone \Subst{N}{Q}{x} 
  \]
  यदि $P \bredone Q$, तो हम कहते हैं कि $P$ \emph{$\beta$-संकुचित होकर}
  $Q$ बनता है। $(\lambd[x][N])Q$ रूप के पद को \emph{रिडेक्स} कहते हैं।
\end{defn}

\begin{prob} \ollabel{prob:def}
  $\beta$-संकुचन की तुल्य आगमनात्मक परिभाषाएँ स्पष्ट रूप से लिखें, जैसा
  हमने \olref[lam][syn][alp]{defn:aconvone} में आबद्ध चर के परिवर्तन के लिए
  किया था।
\end{prob}
  
\begin{defn}[$\beta$-अपचयन, $\bred$] \ollabel{defn:betared}
  \emph{$\beta$-अपचयन} ($\bred$) पदों पर वह सबसे छोटा परावर्ती, संक्रामक
  संबंध है जिसमें $\bredone$ समाहित है। यदि $P \bred Q$, तो हम कहते हैं
  कि $P$, ~$\beta$-अपचयित होकर $Q$ बनता है।
\end{defn}

प्रसंग स्पष्ट होने पर हम $\bredone$ के बदले $\redone$ और $\bred$ के बदले
$\red$ लिखेंगे।

अनौपचारिक रूप से, $M \bred N$ तभी और केवल तभी जब $\beta$-संकुचन के शून्य
या अनेक चरणों द्वारा $M$ को $N$ में बदला जा सके।

\begin{defn}[$\beta$-प्रसामान्य]
जिस पद को आगे $\beta$-संकुचित न किया जा सके वह
\emph{$\beta$-प्रसामान्य} कहलाता है।
\end{defn}

यदि $M \bred N$ और $N$ ~$\beta$-प्रसामान्य है, तो हम $N$ को~$M$ का
\emph{प्रसामान्य रूप} कहते हैं। पूछा जा सकता है कि क्या किसी पद का
प्रसामान्य रूप अद्वितीय है; जैसा हम आगे देखेंगे, उत्तर हाँ है।

कुछ उदाहरणों पर विचार करें।
\begin{enumerate}
\item हमारे पास है
\begin{align*}
(\lambd[x][xxy]) \lambd[z][z] & \redone (\lambd[z][z])(\lambd[z][z]) y \\
& \redone (\lambd[z][z]) y \\
& \redone y
\end{align*}
\item किसी पद को ``सरल करना'' वास्तव में उसे अधिक जटिल बना सकता है:
\begin{align*}
(\lambd[x][xxy])(\lambd[x][xxy]) & \redone (\lambd[x][xxy])(\lambd[x][xxy])y \\
& \redone (\lambd[x][xxy])(\lambd[x][xxy])yy \\
& \redone \dots
\end{align*}
\item यह किसी पद को अपरिवर्तित भी छोड़ सकता है:
\[
(\lambd[x][xx])(\lambd[x][xx]) \redone (\lambd[x][xx])(\lambd[x][xx])
\]
\item कुछ पदों को एक से अधिक प्रकार से अपचयित किया जा सकता है; उदाहरणार्थ,
\[
(\lambd[x][(\lambd[y][yx]) z]) v \redone (\lambd[y][yv]) z
\]
सबसे बाहरी अनुप्रयोग को संकुचित करके; और
\[
(\lambd[x][(\lambd[y][yx]) z]) v \redone (\lambd[x][zx]) v
\]
सबसे भीतरी अनुप्रयोग को संकुचित करके। फिर भी ध्यान दें कि इस स्थिति में
दोनों पद आगे अपचयित होकर एक ही पद~$zv$ बनते हैं।
\end{enumerate}

अंतिम उदाहरण का अंतिम परिणाम संयोग नहीं है; वह लैम्ब्डा कलन के एक गहरे और
महत्त्वपूर्ण गुण को दर्शाता है, जिसे चर्च--रॉसर गुण कहते हैं।

\begin{digress}
  सामान्यतः किसी पद का $\beta$-अपचयन एक से अधिक विधियों से किया जा सकता
  है; इसलिए अनेक अपचयन रणनीतियाँ बनाई गई हैं, जिनमें सबसे सामान्य
  \emph{स्वाभाविक रणनीति} है। स्वाभाविक रणनीति सदा \emph{सबसे बायाँ}
  रिडेक्स संकुचित करती है, जहाँ किसी रिडेक्स की स्थिति पद में उसके आरंभ-बिंदु
  से परिभाषित होती है। स्वाभाविक रणनीति में यह उपयोगी गुण है कि कोई पद किसी
  रणनीति द्वारा प्रसामान्य रूप में तभी और केवल तभी अपचयित हो सकता है जब वह
  स्वाभाविक रणनीति द्वारा प्रसामान्य रूप में अपचयित हो सकता है। आगे अन्यथा
  निर्दिष्ट न होने पर हम स्वाभाविक रणनीति उपयोग करेंगे।
\end{digress}

\begin{defn}[$\beta$-तुल्यता, $\equal$]
  \emph{$\beta$-तुल्यता} ($\equal$) वह संबंध है जो आगमनतः इस प्रकार
  परिभाषित है:
  \begin{enumerate}
  \item $M \equal M$.
  \item यदि $M \equal N$, तो $N \equal M$।
  \item यदि $M \equal N$ और $N \equal O$, तो $M \equal O$।
  \item यदि $M \equal N$, तो $PM \equal PN$।
  \item यदि $M \equal N$, तो $MQ \equal NQ$।
  \item यदि $M \equal N$, तो $\lambd[x][M] \equal \lambd[x][N]$।
  \item $(\lambd[x][N])Q \equal \Subst{N}{Q}{x}$.
  \end{enumerate}
\end{defn}

पहले तीन नियम संबंध को तुल्यता-संबंध बनाते हैं; अगले तीन उसे संगत बनाते
हैं; अंतिम नियम सुनिश्चित करता है कि उसमें $\beta$-संकुचन समाहित है।

अनौपचारिक रूप से, दो पद $\beta$-तुल्य तभी और केवल तभी हैं जब उनमें से एक
को $\beta$-संकुचन अथवा $\beta$-संकुचन के ``प्रतिलोम'' के शून्य या अधिक
चरणों द्वारा दूसरे में बदला जा सके। $\beta$-संकुचन का प्रतिलोम इस प्रकार
परिभाषित है कि $M$ प्रतिलोम-$\beta$-संकुचित होकर $N$ तभी और केवल तभी बनता
है जब $N$, ~$\beta$-संकुचित होकर $M$ बनता है।

ऊपर के नियमों के अतिरिक्त हम संबंध को और नियमों से विस्तारित करेंगे और
विस्तारित तुल्यता-संबंध को $\equal[X]$ से दर्शाएँगे, जहाँ $X$ विस्तारण नियम
है।

\end{document}
